\section{基础模型与通用 AI 的愿景}

\begin{figure}[H]
\centering
\includegraphics[width=0.85\textwidth]{slides-images/slide-77.jpg}
\caption{Foundation Models 催生的强大理念\protect\footnotemark}
\end{figure}
\footnotetext{Slides 第77页。}

在讲座的最后，Amini 将视野拓展到 Foundation Models 的更宏大愿景：

\begin{importantbox}{Foundation Models 的核心愿景}
\begin{itemize}
    \item \textbf{中央推理系统}：生成式基础模型能否为通用 AI 提供一个统一的推理引擎？
    \item \textbf{AI 设计 AI}：能否用 AI 来改进和进化 AI 本身？
    \item \textbf{跨领域生成式 AI}：从图像、生物学、语言到更多领域——既有巨大潜力，也需谨慎对待
    \item \textbf{人工智能与人类智能的关系}：理解两者之间的联系与差异
\end{itemize}
\end{importantbox}

这些问题不仅是技术挑战，更涉及深层的科学和哲学思考。正如 Amini 所强调的：力量与责任并存（power and caution），我们在推动 AI 前沿的同时，必须审慎地考虑其社会影响。

%% =====================================================
